% ========== sections/06_conclusiones.tex ==========
\section{Conclusiones}

\subsection{Respuesta a las preguntas de investigación}

A continuación, se responden las preguntas de investigación planteadas en la introducción, con base en los resultados obtenidos.

\textbf{Pregunta 1: ¿Cuántos grupos naturales de canciones existen en el Top 50 mundial de Spotify?}

Los resultados de ambos algoritmos muestran que la música popular mundial no es homogénea. El algoritmo K-Means identificó 2 grupos naturales, mientras que DBSCAN identificó 8 grupos más 12 puntos de ruido. La diferencia se debe a la naturaleza de cada algoritmo: K-Means tiende a agrupar los datos en pocos clusters de gran tamaño, mientras que DBSCAN, al ser basado en densidad, es capaz de detectar clusters más pequeños y específicos. Considerando ambos enfoques, se puede concluir que existen al menos 2 grandes perfiles musicales dominantes en el Top 50, con múltiples subperfiles de nicho.

\textbf{Pregunta 2: ¿Es posible asignar nombres descriptivos a los clusters basándose en sus valores promedio de danceability, energy y valence?}

Sí, es posible. Utilizando exclusivamente las variables danceability, energy y valence, se distinguen claramente los diferentes perfiles musicales. Para K-Means se identificaron dos perfiles: \enquote{Bailable, energética y positiva} (alta danceability, alta energy, alta valence) y \enquote{Mixta / equilibrada} (valores moderados en todas las dimensiones). Para DBSCAN se identificaron perfiles como \enquote{Popular comercial}, \enquote{Instrumental introspectivo}, \enquote{Ambiental extremo}, \enquote{Balada lenta}, \enquote{Rap hablado}, \enquote{Folk alegre}, \enquote{Estudio y concentración} e \enquote{Infantil y lullabies}.

\textbf{Pregunta 3: ¿Qué combinación de valores caracteriza a cada perfil musical?}

Las combinaciones características para K-Means son:

\begin{itemize}
    \item \textbf{Cluster 1 (Bailable, energética y positiva):} danceability 0.699, energy 0.743, valence 0.611.
    \item \textbf{Cluster 0 (Mixta / equilibrada):} danceability 0.592, energy 0.491, valence 0.382.
\end{itemize}

Para DBSCAN, las combinaciones más representativas son:

\begin{itemize}
    \item \textbf{Popular comercial:} danceability 0.678, energy 0.652, valence 0.549.
    \item \textbf{Folk alegre:} danceability 0.531, energy 0.038, valence 0.800.
    \item \textbf{Balada lenta:} danceability 0.094, energy 0.002, valence 0.068.
\end{itemize}

\textbf{Pregunta 4: ¿Qué algoritmo de clustering produce los mejores resultados en términos de coeficiente de Silhouette?}

El algoritmo DBSCAN obtuvo un coeficiente de Silhouette de 0.516, superior al de K-Means (0.210). Esto indica que DBSCAN produce clusters mejor definidos y más separados entre sí. Sin embargo, K-Means tiene la ventaja de no generar puntos de ruido (0\% de datos no asignados) y ser más simple de interpretar. La elección del algoritmo depende del objetivo: DBSCAN es preferible si se busca identificar clusters de nicho y detectar datos atípicos, mientras que K-Means es más adecuado para una segmentación rápida y general del mercado musical.

\subsection{Aportes y hallazgos principales}

Los principales aportes de este trabajo son los siguientes:

\begin{enumerate}
    \item Se ha demostrado que es posible segmentar los grandes éxitos de Spotify en perfiles acústicos diferenciados, identificando al menos 2 grandes grupos y múltiples subgrupos de nicho.
    
    \item Se ha propuesto un sistema de nomenclatura semántica para los clusters (Bailable, energética y positiva; Popular comercial; Folk alegre; Instrumental introspectivo, etc.), facilitando la comunicación de resultados a audiencias no técnicas.
    
    \item Se ha establecido una comparativa cuantitativa entre K-Means y DBSCAN, evidenciando que DBSCAN ofrece mejor separación de clusters (Silhouette 0.516 vs 0.210), mientras que K-Means es más rápido de ejecutar y más simple de interpretar.
    
    \item Se han formulado caracterizaciones acústicas para la industria musical, especificando qué combinaciones de danceability, energy y valence caracterizan a cada perfil de éxito.
    
    \item Se ha diferenciado explícitamente este trabajo de análisis públicos existentes en Kaggle, destacando el enfoque en clustering no supervisado, la comparación de dos algoritmos (K-Means y DBSCAN) y la asignación de nombres interpretables a los clusters.
\end{enumerate}

\subsection{Limitaciones y trabajo futuro}

\textbf{Limitaciones del estudio:}

\begin{itemize}
    \item El análisis se limita a canciones que ya alcanzaron el Top 50, introduciendo un sesgo de selección. No se analizan canciones que no lograron éxito.
    
    \item El coeficiente de Silhouette de DBSCAN (0.516) indica una calidad de clustering buena pero no excelente. Existe cierta superposición entre clusters.
    
    \item El muestreo a 100.000 canciones, aunque necesario por costo computacional, podría haber excluido algunos patrones presentes en el dataset completo de 2.1 millones de registros.
    
    \item No se incluyeron variables contextuales como género musical, año de lanzamiento o país de origen, que podrían enriquecer la segmentación.
\end{itemize}

\textbf{Trabajo futuro:}

\begin{itemize}
    \item Aplicar técnicas de reducción de dimensionalidad alternativas (t-SNE, UMAP) para visualizar mejor la separación entre clusters.
    
    \item Incorporar variables temporales para analizar cómo evolucionan los perfiles acústicos de los éxitos a lo largo del tiempo.
    
    \item Entrenar un clasificador supervisado que, dadas las características acústicas de una canción nueva, prediga su perfil cluster.
    
    \item Repetir el análisis por país o región para identificar diferencias culturales en los perfiles de éxito musical.
\end{itemize}

\subsection{Aprendizajes del proceso}

La realización de este proyecto ha permitido consolidar aprendizajes fundamentales en el ciclo de vida de un proyecto de Ciencia de Datos:

\begin{itemize}
    \item \textbf{Importancia del preprocesamiento:} Se confirmó que el preprocesamiento y la limpieza de datos consumen una parte significativa del tiempo del proyecto, siendo una fase crítica para obtener resultados confiables.
    
    \item \textbf{No hay algoritmo universalmente mejor:} La elección del algoritmo de clustering depende del objetivo del análisis (segmentación general vs. detección de nichos) y de las características del dataset.
    
    \item \textbf{Valor de la interpretabilidad:} La asignación de nombres descriptivos a los clusters es tan importante como los resultados numéricos, especialmente para comunicar hallazgos a audiencias no técnicas.
    
    \item \textbf{Trabajo colaborativo:} Se confirmó la importancia del trabajo en equipo con control de versiones (GitHub) y de una estructura modular del proyecto (datos, notebooks, documentación) para facilitar la integración de contribuciones individuales.
    
    \item \textbf{Iteración constante:} El proceso de análisis de datos es inherentemente iterativo; los resultados iniciales guían nuevas preguntas y refinamientos del modelo.
\end{itemize}